Malgré ces difficultés, notre montagne va faire parler d'elle grâce
au capitaine Pujol de la Tapie qui va se mettre à la tête d'une troupe
catholique et se couvrir de quelque gloire... Les Pujol de la Tapie
(dont les Saluste descendent), sont une vieille famille du pays.
Un acte notarié de 1301 signale que Jean et Bertrand Pujol de la
Tapie reçoivent de Gaillarde de Roquefort des bois et « boscals »
sis entre Paulin et Massals. Mais c'est surtout au XVII\textsuperscript{ème}
siècle que cette famille fait parler d'elle. On trouve un Jean Pujol
de la Tapie, notaire à Plaisance, un Pierre chapelain de Cambon d'Albi,
un Antoine « procureur général » au siège royal de Plaisance,
lieutenant de la justice du vicomté de Paulin et de la première
compagnie de bourgeoisie du dit lieu, habitant la Tapie. Mais avec
la guerre de Rohan, c'est Jehan de la Tapie qui a la vedette. On le
prétend le « meilleur pétardier de France et de Navarre, n'ayant
jamais chargé et posé pétard qui n'eut fait effet ». On le surnomme
« l'artillogat » et il fait trembler les huguenots. On eut juger\footnote{???} de
son prestige si l'on pense qu'il lui suffisait de sonner du cor pour
voir appliquer à la Tapie, en une demi-heure, une quinzaine de
partisans. Une journée lui suffisait pour en avoir une centaine à
sa disposition en son castel de la Tapie. Avec ses montagnards,
il s'était joint à l'armée du grand Condé... Et tandis que Mr de
la Tapie était au loin, des brigands surnommés « les cadets de
Massals » venaient commettre des dégâts à Alban et ailleurs.
Pour en avoir raison, le prévôt de Castres dut venir les assiéger
dans le château de Massals devenu leur repaire. C'était en
novembre 1623.

Les bagarres reprendront sous peu. Dans l'intervalle, en mai 1624,
les biens de l'église de Carménel sont arrentés à Mas de Massals
pour la somme de 550 livres.

Dans un camp comme dans un autre, le racisme religieux continue
à sévir. Les braves gens en pâtissent. Lacaune qui est dans le
camp des réformés en fait voir de cruelles aux catholiques.
Les registres de 1622 signalent que les maîtres de la cité obligent
les « paysans papistes à se convertir » sinon on enlève la toiture
de leurs maisons et on met en prison ces catholiques irréductibles.